\begingroup
\setlength{\LTcapwidth}{\textwidth}
\begin{longtable}{p{0.23\textwidth}p{0.27\textwidth}p{0.37\textwidth}}
\caption{Definisi dan rasionalisasi fitur modul prediksi.}
\label{tab:tab5.5_definisi_fitur_modul_prediksi} \\
\toprule
\textbf{Fitur} & \textbf{Transformasi} & \textbf{Rasionalisasi penggunaan} \\
\midrule
\endfirsthead

\multicolumn{3}{@{}l}{\small\emph{Tabel \thetable{} (lanjutan)}} \\
\toprule
\textbf{Fitur} & \textbf{Transformasi} & \textbf{Rasionalisasi penggunaan} \\
\midrule
\endhead

\bottomrule
\endlastfoot

\texttt{glucose} & Nilai glukosa pada waktu pengamatan & Menyediakan keadaan kadar glukosa yang diketahui \\

\texttt{glucose\_delta} & Selisih antar-pengamatan & Menangkap arah perubahan kadar glukosa \\

\texttt{iob} & Peluruhan insulin berbasis waktu & Mempertahankan efek residual insulin \\

\texttt{cob} & Peluruhan karbohidrat berbasis waktu & Mempertahankan efek residual karbohidrat \\

\texttt{activity} & Nilai aktivitas yang tersedia & Menambahkan konteks aktivitas fisik \\

\texttt{hour\_sin} & Sinus waktu, periode 24 jam & Merepresentasikan pola waktu secara siklis \\

\texttt{hour\_cos} & Cosinus waktu, periode 24 jam & Melengkapi representasi siklis waktu \\

\texttt{glucose\_rate} & Perubahan glukosa per waktu & Menangkap kecepatan perubahan glukosa \\

\texttt{time\_since\_\allowbreak prev\_glucose} & Selang waktu aktual & Mempertahankan irregularitas pengamatan \\
\end{longtable}
\endgroup
